\documentclass[10pt,twocolumn]{article}
\usepackage[T1]{fontenc}
\usepackage[utf8]{inputenc}
\usepackage{lmodern}
\usepackage[margin=0.7in]{geometry}
\usepackage{graphicx}
\usepackage{microtype}
\usepackage{caption}

\title{Font Bigger Demo Paper}
\author{VS Code Prototype}
\date{}

\begin{document}
\maketitle

\section{Why This Exists}

This minimal paper demonstrates a tight loop for tuning Python-generated figures in the context that matters: the final PDF layout.

\section{Figure In Context}

The figure below is intentionally included at column width. The default style file uses a larger font so the plot stays readable after typesetting.

\begin{figure}[t]
  \centering
  \includegraphics[width=\columnwidth]{figures/output.pdf}
  \caption{A matplotlib figure generated from a JSON style file.}
  \label{fig:demo}
\end{figure}

Small parameter tweaks are usually obvious only after the figure is embedded in the paper. This prototype shortens that feedback loop.

\end{document}
